\clearpage\begingroup\SourceLeaf{83b}\HistoricalHead{78}{L}{PHYSIK. BRIEFWECHSEL.}
\begin{center}[6.]\end{center}
\begin{center}\textsc{Gauss} an \textsc{Gerling}.\quad Göttingen, 20. Junius 1832.\end{center}
\noindent . . . . . . . . . . . . . . . . . . . . .\par
\indent Meine Zurückführung der Intensität auf absolute Einheit, wozu ich
schon mehrere, obwohl erst als vorläufig anzusehende Versuche gemacht habe,
gelingen ganz unvergleichlich. Aber das von \textsc{Moser} und \textsc{Rieser}[\textsuperscript{*})]\begingroup\renewcommand{\thefootnote}{}
\footnotetext{[\textsuperscript{*}) Siehe die Fußnoten auf S. 77 und 75.]}
\endgroup\ aus der
Beobachtung in Berlin berechnete Resultat ist nur \(\frac{1}{5}\) des meinigen, also ganz
\clearpage\SourceLeaf{84a}\HistoricalHead{79}{R}{GAUSS AN OLBERS.}
\noindent unbrauchbar (mein Resultat bestätigt sich auch durch Versuche an Nadeln
von den verschiedensten Dimensionen, obwohl kleine Nadeln wenig Genauigkeit
geben können). Jener enorme Fehler hat übrigens seinen Grund hauptsächlich
in einer ganz unzulässigen Berechnungsweise: nach richtigen Principien
finden sich, so gut es geht, Resultate, die wenigstens Annäherungen
sind und sogar mein Resultat zwischen sich haben. . . . . . . . . . . . . . . .
\clearpage\endgroup
